% GENERATED FILE. Edit canonical lessons, then run npm run generate:books.
% canonical-source-hash: fnv1a64:7257bffcd217f84e
% canonical-lessons: FR-C495-dispute, FR-C495-brouiller, FR-C495-tromper, FR-C495-trahir, FR-C495-trahison

\chapter{Disputes et trahisons}
\label{ch:fr-disputes-et-trahisons}
\hlchaptermodality{495}

\begin{chapteropening}
\textbf{By the end of this chapter:} \emph{I can talk about quarrels, betrayal and falling out.}

Complete the last lesson of chapter 495: I can talk about quarrels, betrayal and falling out.
\end{chapteropening}

\section[la dispute]{la dispute — quarrel, argument}

\label{lesson:FR-C495-dispute}



\begin{quote}
Before the new one: say the French for to disobey, then the French for day nursery.
\end{quote}

\subsection*{You'll want to know: la dispute}
\textbf{la dispute} — \textquotedblleft{}quarrel, argument\textquotedblright{}.

\begin{quote}
Ils ont eu une dispute à propos de l'argent.
\end{quote}

\begin{quote}
\emph{They had an argument about money.}
\end{quote}

A real row, not a debate.

\begin{cousinweb}
\emph{la dispute} is built on the verb \emph{disputer}, from Latin \emph{disputāre} \textquotedblleft{}to discuss, to argue\textquotedblright{} (EtymDB), built from \emph{dis-} \textquotedblleft{}apart\textquotedblright{} and \emph{putāre} \textquotedblleft{}to reckon\textquotedblright{} (Lewis \& Short).
\end{cousinweb}

\begin{grammarlens}[title={What you've built}]
One more word for telling, explaining and describing at B1.
\end{grammarlens}

\subsection*{Your turn}
\begin{itemize}
\raggedright
  \item \emph{Say it:} \emph{la dispute}
  \item \emph{Say it:} \emph{la dispute}, once more
  \item \emph{Say it:} a sentence of your own with \emph{la dispute}
  \item \emph{Recall:} say the French for to disobey, then the French for day nursery, then say \emph{la dispute} again
\end{itemize}

\begin{tcolorbox}[breakable,colback=teal!4,colframe=teal!35!black,title={Before you move on}]
What does la dispute mean? (\textquotedblleft{}Quarrel, argument\textquotedblright{}.) Say it once more.
\end{tcolorbox}
\section[brouiller]{brouiller — to mix up, to blur; (se brouiller) to fall out}

\label{lesson:FR-C495-brouiller}



\begin{quote}
Before the new one: say the French for day nursery, then the French for quarrel.
\end{quote}

\subsection*{You'll want to know: brouiller}
\textbf{brouiller} — \textquotedblleft{}to mix up, to blur; (se brouiller) to fall out\textquotedblright{}.

\begin{quote}
Les deux frères se sont brouillés à cause d'une maison.
\end{quote}

\begin{quote}
\emph{The two brothers fell out over a house.}
\end{quote}

\emph{Se brouiller avec quelqu'un} — to fall out with someone. \emph{Des œufs brouillés} — scrambled eggs.

\begin{cousinweb}
\emph{brouiller} goes back to Old French \emph{brouiller}, from Vulgar Latin \emph{brodiculāre} \textquotedblleft{}to jumble together\textquotedblright{} (EtymDB), which is usually traced to a Germanic word for broth.
\end{cousinweb}

\begin{grammarlens}[title={What you've built}]
One more word for telling, explaining and describing at B1.
\end{grammarlens}

\subsection*{Your turn}
\begin{itemize}
\raggedright
  \item \emph{Say it:} \emph{brouiller}
  \item \emph{Say it:} \emph{brouiller}, once more
  \item \emph{Say it:} a sentence of your own with \emph{brouiller}
  \item \emph{Recall:} say the French for day nursery, then the French for quarrel, then say \emph{brouiller} again
\end{itemize}

\begin{tcolorbox}[breakable,colback=teal!4,colframe=teal!35!black,title={Before you move on}]
What does brouiller mean? (\textquotedblleft{}To mix up, to blur; (se brouiller) to fall out\textquotedblright{}.) Say it once more.
\end{tcolorbox}
\section[tromper]{tromper — to deceive; to be unfaithful to}

\label{lesson:FR-C495-tromper}



\begin{quote}
Before the new one: say the French for quarrel, then the French for to mix up.
\end{quote}

\subsection*{You'll want to know: tromper}
\textbf{tromper} — \textquotedblleft{}to deceive; to be unfaithful to\textquotedblright{}.

\begin{quote}
Elle a appris que son compagnon la trompait.
\end{quote}

\begin{quote}
\emph{She found out that her partner was cheating on her.}
\end{quote}

\emph{Se tromper} — to make a mistake, a very common sense.

\begin{cousinweb}
\emph{tromper} goes back to Old French \emph{tromper} (EtymDB). The origin is not settled; a link with \emph{trompe} \textquotedblleft{}horn\textquotedblright{} is often proposed.
\end{cousinweb}

\begin{grammarlens}[title={What you've built}]
One more word for telling, explaining and describing at B1.
\end{grammarlens}

\subsection*{Your turn}
\begin{itemize}
\raggedright
  \item \emph{Say it:} \emph{tromper}
  \item \emph{Say it:} \emph{tromper}, once more
  \item \emph{Say it:} a sentence of your own with \emph{tromper}
  \item \emph{Recall:} say the French for quarrel, then the French for to mix up, then say \emph{tromper} again
\end{itemize}

\begin{tcolorbox}[breakable,colback=teal!4,colframe=teal!35!black,title={Before you move on}]
What does tromper mean? (\textquotedblleft{}To deceive; to be unfaithful to\textquotedblright{}.) Say it once more.
\end{tcolorbox}
\section[trahir]{trahir — to betray}

\label{lesson:FR-C495-trahir}



\begin{quote}
Before the new one: say the French for to mix up, then the French for to deceive.
\end{quote}

\subsection*{You'll want to know: trahir}
\textbf{trahir} — \textquotedblleft{}to betray\textquotedblright{}.

\begin{quote}
Elle a trahi le secret de son amie.
\end{quote}

\begin{quote}
\emph{She betrayed her friend's secret.}
\end{quote}

A regular \emph{-ir} verb: \emph{nous trahissons}.

\begin{cousinweb}
\emph{trahir} is inherited from Latin \emph{trādere} \textquotedblleft{}to hand over\textquotedblright{} (Lewis \& Short), built from \emph{trāns} and \emph{dare} \textquotedblleft{}to give\textquotedblright{}.
\end{cousinweb}

\begin{grammarlens}[title={What you've built}]
One more word for telling, explaining and describing at B1.
\end{grammarlens}

\subsection*{Your turn}
\begin{itemize}
\raggedright
  \item \emph{Say it:} \emph{trahir}
  \item \emph{Say it:} \emph{trahir}, once more
  \item \emph{Say it:} a sentence of your own with \emph{trahir}
  \item \emph{Recall:} say the French for to mix up, then the French for to deceive, then say \emph{trahir} again
\end{itemize}

\begin{tcolorbox}[breakable,colback=teal!4,colframe=teal!35!black,title={Before you move on}]
What does trahir mean? (\textquotedblleft{}To betray\textquotedblright{}.) Say it once more.
\end{tcolorbox}
\section[la trahison]{la trahison — betrayal; treason}

\label{lesson:FR-C495-trahison}



\begin{quote}
Before the new one: say the French for to deceive, then the French for to betray.
\end{quote}

\subsection*{You'll want to know: la trahison}
\textbf{la trahison} — \textquotedblleft{}betrayal; treason\textquotedblright{}.

\begin{quote}
Pour lui, ce mensonge est une trahison.
\end{quote}

\begin{quote}
\emph{For him, this lie is a betrayal.}
\end{quote}

Built on \emph{trahir}.

\begin{cousinweb}
\emph{la trahison} comes from Old French \emph{traison} (EtymDB), built on \emph{trair} \textquotedblleft{}to betray\textquotedblright{}, from Latin \emph{trādere} \textquotedblleft{}to hand over\textquotedblright{}, built on \emph{dare} (Lewis \& Short).
\end{cousinweb}

\begin{grammarlens}[title={What you've built}]
One more word for telling, explaining and describing at B1.
\end{grammarlens}

\subsection*{Your turn}
\begin{itemize}
\raggedright
  \item \emph{Say it:} \emph{la trahison}
  \item \emph{Say it:} \emph{la trahison}, once more
  \item \emph{Say it:} a sentence of your own with \emph{la trahison}
  \item \emph{Recall:} say the French for to deceive, then the French for to betray, then say \emph{la trahison} again
\end{itemize}

\begin{tcolorbox}[breakable,colback=teal!4,colframe=teal!35!black,title={Before you move on}]
What does la trahison mean? (\textquotedblleft{}Betrayal; treason\textquotedblright{}.) Say it once more.
\end{tcolorbox}
